\section{二分查找、边界和答案单调性}

二分查找最难的不是中点公式，而是把问题改写成边界问题。搜索插入位置找的是第一个大于等于目标的位置，查找首尾区间找的是两个边界，旋转数组找的是有序半边，答案二分找的是可行和不可行的分界。只要边界语义清楚，循环条件、左右更新和返回值就会自然统一；如果边界语义模糊，就会陷入各种加一减一的补丁。

推荐固定左闭右开区间 [left,right)。初始 right 等于候选数量，循环条件是 left 小于 right，中点落在合法候选中。当谓词在 mid 处为假，说明答案在 mid 右侧，于是 left 等于 mid 加一；当谓词为真，说明 mid 可能就是答案，于是 right 等于 mid。循环结束时 left 等于 right，是第一个满足谓词的位置。这个模板能覆盖 lower\_bound、答案最小可行值和许多边界题。

答案二分的关键是谓词单调。爱吃香蕉中，速度越快越容易完成；船运包裹中，容量越大需要天数越少；分割数组中，允许的最大段和越大，所需段数越少；最小体力路径中，体力阈值越大，可走边越多。这些问题的原数组未必有序，但答案空间有序。若检查函数内部使用了不正确的贪心，外层二分再漂亮也没有意义。

写检查函数时要先说明贪心为什么正确。船运包裹按顺序装包，当前天能放就放，放不下再开新天；因为包裹顺序不能改变，提前开新天只会让剩余天数更紧，不会更优。分割数组最大值也是同理：在段和不超过上限的前提下尽量延长当前段，可以最小化段数。检查函数本身常常就是一个贪心证明题。

二分的溢出问题不能忽略。中点用 left 加 (right-left)/2，平方比较用 mid <= x/mid，求总耗时或总容量时用 \code{std::int64_t}。很多题的输入范围故意接近 int 上限，代码逻辑正确但乘法溢出会得到完全错误的判断。复杂度分析也要写成 O(n log M)，其中 M 是答案范围，不要笼统写 O(log n)。

旋转数组和峰值题展示了另一类二分：不是直接找单调谓词，而是利用局部比较保证某一侧存在答案。寻找峰值中，如果 nums[mid] 小于 nums[mid+1]，右侧一定有峰；否则左侧含 mid 一定有峰。这个证明依赖边界可视为负无穷和相邻不相等的条件。若题目条件改变，例如允许大量重复，就要重新证明是否还能排除一半。

实验建议：把同一道 lower\_bound 用左闭右闭、左闭右开两种写法实现，再用长度从零到五的所有数组暴力枚举测试。答案二分实验中，先写一个很慢但显然正确的枚举版本，再和二分版本对比。对于每个失败用例，记录 left、mid、right 和谓词值，直到能指出哪条不变量被破坏。

\begin{principlebox}{本节复盘}
阅读本节后，请回到相邻章节的例题，把暴力瓶颈、优化依据、状态不变量、C++ 容器成本和边界测试重新写一遍。能把同一思想迁移到变体题，才算真正掌握。
\end{principlebox}
